\documentclass[11pt]{article}
\newcommand{\courseshort}{Data Structures (Make-up) · 010153103}
\newcommand{\lecturetitle}{L7: Priority Queues \& Heaps}
\usepackage{lecture_style}

\begin{document}
\lectureheader{Lecture 7: Priority Queues \& Heaps}{Pre-class brief}{Goodrich Ch.~9}

\begin{objbox}
\begin{itemize}[leftmargin=*]
  \item State the priority-queue ADT and its main operations.
  \item Define the heap-order property and the structural property of a heap.
  \item Implement a binary heap with an array using \emph{up-heap} and \emph{down-heap}.
  \item Outline the heap-sort algorithm and its $O(n\log n)$ running time.
\end{itemize}
\end{objbox}

\section*{1. Priority Queue (PQ) ADT}

\begin{itemize}[leftmargin=*]
  \item \texttt{insert(k, v)} --- insert an entry with key (priority) $k$.
  \item \texttt{min()} --- return the entry with smallest key.
  \item \texttt{removeMin()} --- remove and return that entry.
\end{itemize}

\section*{2. Binary Heap}

A \textbf{heap} is a complete binary tree where every parent's key is $\leq$ its children's keys (\emph{min-heap}).

\textbf{Array representation.} Store the tree level by level in an array:
\begin{itemize}[leftmargin=*]
  \item Root at index 0 (or 1 --- pick a convention).
  \item Children of index $i$: $2i+1$ and $2i+2$. Parent: $(i-1)/2$.
\end{itemize}

\begin{center}
\begin{tikzpicture}[
  node/.style={circle, draw, minimum size=0.7cm, font=\small\sffamily}
]
  \node[node] (a) {4};
  \node[node, below left=0.8cm and 1.2cm of a] (b) {7};
  \node[node, below right=0.8cm and 1.2cm of a] (c) {6};
  \node[node, below left=0.7cm and 0.3cm of b] (d) {9};
  \node[node, below right=0.7cm and 0.3cm of b] (e) {8};
  \draw (a)--(b); \draw (a)--(c); \draw (b)--(d); \draw (b)--(e);
\end{tikzpicture}\\[2pt]
{\small array: $[4, 7, 6, 9, 8]$\quad indices: $[0, 1, 2, 3, 4]$}
\end{center}

\section*{3. Insert (up-heap)}

Append at the end, then \emph{swap with parent} while smaller. Cost: $O(\log n)$.

\begin{lstlisting}
public void insert(int k) {
    h[size++] = k;
    int i = size - 1;
    while (i > 0 && h[(i-1)/2] > h[i]) {
        swap(i, (i-1)/2);
        i = (i-1)/2;
    }
}
\end{lstlisting}

\section*{4. Remove-min (down-heap)}

Replace root with last element, shrink, then swap with the smaller child while needed. Cost: $O(\log n)$.

\begin{lstlisting}
public int removeMin() {
    int min = h[0];
    h[0] = h[--size];
    int i = 0;
    while (2*i + 1 < size) {
        int c = 2*i + 1;
        if (c+1 < size && h[c+1] < h[c]) c++;
        if (h[i] <= h[c]) break;
        swap(i, c); i = c;
    }
    return min;
}
\end{lstlisting}

\section*{5. Heap Sort (preview)}

Insert all $n$ elements ($O(n\log n)$), then call \texttt{removeMin} $n$ times ($O(n\log n)$). Total: $O(n\log n)$. The sequence of removed elements is sorted.

\begin{exbox}{Pre-Class Exercises (do BEFORE the lecture)}
\textbf{Exercise 7.1 --- Trace insertions.} Start with an empty min-heap. Insert in order: \verb|7, 4, 9, 1, 8, 5|. Draw the heap and the array after each insertion.

\textbf{Exercise 7.2 --- Trace removals.} From the heap \texttt{[2, 5, 4, 8, 7, 9, 6]} (a valid min-heap), perform \texttt{removeMin} three times and show the heap after each step.

\textbf{Exercise 7.3 --- Heap or not?} For each array, decide whether it is a min-heap. If not, mark the violating parent--child pair.
\begin{enumerate}[label=(\alph*)]
  \item $[1, 3, 2, 5, 4, 9, 8]$
  \item $[2, 4, 6, 8, 5, 9]$
  \item $[1, 6, 5, 7, 8, 9, 10]$
\end{enumerate}

\textbf{Exercise 7.4 --- Implement.} Code the \verb|insert| method for a max-heap (parent's key $\geq$ children's keys).

\textbf{Exercise 7.5 --- Why a heap, not a sorted array?} An entrepreneur says: ``I'll just keep priorities in a sorted array; \texttt{min} is then $O(1)$!'' Compare insert costs and explain when each design wins.

\textbf{Exercise 7.6 --- $k$-th smallest.} Given an array of $n$ numbers, describe an algorithm using a heap that finds the $k$-th smallest in $O(n\log k)$. (Hint: a max-heap of size $k$.)
\end{exbox}

\begin{disbox}
\begin{itemize}[leftmargin=*]
  \item Why is the heap stored in an array rather than as a linked tree? Memory and cache reasons.
  \item Discuss Exercise 7.6 --- this is a classic streaming problem (top-$k$).
  \item What does \texttt{PriorityQueue<E>} in Java actually use? When is it the right choice?
\end{itemize}
\end{disbox}

\end{document}
